\documentclass[10pt,a4paper]{amsart}
\usepackage[margin=19mm]{geometry}
\usepackage{amsmath,amssymb,mathtools}
\usepackage[scr=rsfs,cal=euler]{mathalfa}
\usepackage{xfrac}
\usepackage[unicode,hidelinks,bookmarks=false]{hyperref}
\newcommand{\ie}{i.e.\ }
\newcommand{\cf}{cf.\ }
\newcommand{\resp}{respectively\ }
\newcommand{\Subsection}{\subsection}
\providecommand{\nobreakdash}{\nobreak\hbox{-}}
\setlength{\emergencystretch}{2em}
\multlinegap0pt
\usepackage{siunitx}
\usepackage{siunitx}
\usepackage{siunitx}
\usepackage{siunitx}
\usepackage{algorithm}\usepackage{algpseudocode}
\usepackage{xcolor}
\usepackage{tikz}
\usepackage{bm}
\usepackage{mathtools}
\usepackage{amssymb}
\usepackage{float}
\usepackage{siunitx}
\newtheorem{theorem}{Theorem}
\newtheorem{lemma}{Lemma}
\usepackage{cleveref}
\crefname{theorem}{Theorem}{Theorems}
\crefname{lemma}{Lemma}{Lemmas}
\renewcommand{\P}{\mathbb{P}}
\DeclareMathOperator{\R}{\mathbb{R}}
\newcommand{\diff}{\, \mathrm{d}}
\title{On the Effectiveness of Classical Regression Methods for Optimal Switching Problems}
\author{Martin Andersson, Benny Avelin, Marcus Olofsson}
\date{Source: arXiv:2506.15436v4 (CC BY 4.0)}
\begin{document}
\maketitle

\noindent\emph{Source and licence.} On the Effectiveness of Classical Regression Methods for Optimal Switching Problems. Version arXiv:2506.15436v4 (CC BY 4.0), distributed by arXiv under the Creative Commons Attribution 4.0 International licence, http://creativecommons.org/licenses/by/4.0/, as recorded in the arXiv licence metadata for this exact version and shown on the abstract page https://arxiv.org/abs/2506.15436v4. Full licence notice: \texttt{licenses/arxiv-2506.15436-CC-BY-4.0.txt}.\par\smallskip

\noindent\textit{Editorial scope.} Counted unit: the article's lemma lemma:gj:lipschitz (source0.tex lines 1741--1747). The article's complete original proof of it is delivered with it (source0.tex lines 1748--1765, one \texttt{proof} environment of 18 lines). No mathematics is written, repaired or re-derived: the statement and the proof are the article's own lines, in the article's own notation, and no retained line is substituted or re-flowed.  Retained uncounted as the source-local context the proof needs: lines 292--294; lines 510--516; lines 539--541; lines 554--570.  Nothing was omitted from any retained span, so the package carries no omission record.  No retained line contains a citation, so the wrapper carries no bibliography and no bibliography entry.  No retained line contains a figure environment, an image-inclusion command or drawing code, so no image is delivered and the package declares no asset dependency.  Editorial formatting only, in addition: an AMS \texttt{amsart} preamble that restates the article's own declarations and macros used by the retained lines, namely the environments \texttt{theorem}, \texttt{lemma} on separate counters, together with the standard packages those lines need.  The retained environments print in this document as Theorem 1, Lemma 1 in order of appearance, whereas the article prints the same items with the numbering of its own full document.  The complete native source of the article as distributed in the arXiv e-print for this exact version is bundled unchanged as \texttt{sources/180088/source0.tex}.
\begin{equation} \label{eq:gj}
  g_j(t_n,x) = \int_{\R^d}{\widehat V}_j(t_{n+1},y)\rho(x,t_n;y,t_n+\Delta t)\,dy,
\end{equation}

\begin{equation} \label{eq:transition:bounds}
  \begin{split}
    |\rho(x,t;y,t+\Delta t)|          & \leq \frac{1}{\lambda_1} \exp(-\lambda_1 \|x-y\|^2),
    \\
    |\nabla_x \rho(x,t;y,t+\Delta t)| & \leq \frac{1}{\lambda_2} \exp(-\lambda_2 \|x-y\|^2).
  \end{split}
\end{equation}

\begin{equation} \label{eq:truncation}
  \widehat{g}_j(t_n,x) = \mathrm{sign}(\widetilde{g}_j(t_n,x)) \cdot \min\{|\widetilde{g}_j(t_n,x)|, 2 C_0(\|x\|+1)\},
\end{equation}

\begin{theorem} \label{thm:1}
  Let \cref{eq:transition:bounds} hold and assume that $|\widehat{V}_i(t_{n+1},x)| \leq L (\|x\|+1)$ for all $i = 1,\dots,D$ and some $L > 0$.
  Then for any fixed bounded set $Q \subset \R^d$ there exist positive constants
  $C_0$ (used in \cref{eq:truncation}), $C_1$, $C_2$ and $\sigma$, depending only on $d$, $L$, $Q$, $\lambda_1$ and $\lambda_2$, such that for any $\delta > 0$ it holds
  \begin{multline*}
    \P(\sup_{x \in Q} |\widehat{V}_i(t_n,x) - \widetilde{V}_i(t_n,x)| \geq \delta) \leq
    2D{\binom{M}{k}} \exp \left (-\frac{k m_Y\delta^2}{18 \sigma^2} \right )
    \\
    + D \cdot
    \begin{cases}
      C_1 \frac{|Q|}{\delta^d} \exp \left (-\frac{M \left (\frac{M-k}{M} - \sup_Q p_{\delta/(4 C_2)}^c \right )_+^2}{2} \right )            & \text{if } \delta/C_2 \leq 1,
      \\
      C_1 \frac{|Q|}{\delta^{d/2}} \exp \left (-\frac{M \left (\frac{M-k}{M} - \sup_Q p_{\sqrt{\delta/(4 C_2)}}^c \right )_+^2}{2} \right ) & \text{if } \delta/C_2 > 1,
    \end{cases}
  \end{multline*}
  where $D$ is the number of modes, $p_{r}^c = \P(X_{t_n} \not \in B(x,r))$, $B(x,r)$ is the Euclidean ball of radius $r$ centered at $x \in \R^d$, and $|Q|$ is the volume of the set $Q$.
\end{theorem}

\begin{lemma} \label{lemma:gj:lipschitz}
  Consider the function $g_j(t_{n},x)$, as defined in \cref{eq:gj}, and assume the conditions as in \cref{thm:1}.
  Then there exists a constant $C_0(d,L,\lambda_1,\lambda_2)$ such that
  \begin{equation*}
    |g_j(t_{n},x)| + |\nabla_x g_j(t_{n},x)| \leq C_0(\|x\|+1).
  \end{equation*}
\end{lemma}

\begin{proof}
  We can write
  \begin{align*}
    \nabla_x g_j(t_{n},x) & = \int \nabla_x \rho(x,t_n;y,t_{n+1}) \widehat{V}_j(t_{n+1},y) \diff y.
  \end{align*}
  From the growth bounds on $\widehat{V}_j$ together with \cref{eq:transition:bounds} and a change of variables we get that there exist constants $C$ (not necessarily the same at every line) and $\lambda > 0$ such that
  \begin{align*}
    \left | \nabla_x g_j(t_{n},x) \right |
     & \leq
    C \int \|\nabla_x \rho(x,t_n;y,t_{n+1})\| (\|y\|+1) \diff y
    \\
     & \leq
    C \int (\|z+x\|+1) \exp(-\lambda\|z\|^2/\Delta t) \diff z
    \leq
    C (\|x\|+1).
  \end{align*}
  A similar calculation shows that $|g_j(t_{n},x)| \leq C(\|x\|+1)$.
\end{proof}

\end{document}
